\documentclass[11pt]{amsart}
\usepackage[margin=1in]{geometry}
\usepackage{amssymb,amsmath,amsthm,mathtools}
\usepackage{booktabs,array,tabularx}
\usepackage[hidelinks]{hyperref}
\hypersetup{pdftitle={Two-Port Substitution and the Growth of the Smallest Current in Planar Networks},pdfauthor={Filip Majchrzak}}
\IfFileExists{microtype.sty}{\usepackage{microtype}}{}
\emergencystretch=3em

\theoremstyle{plain}
\newtheorem{theorem}{Theorem}
\newtheorem{proposition}[theorem]{Proposition}
\newtheorem{lemma}[theorem]{Lemma}
\newtheorem{corollary}[theorem]{Corollary}
\theoremstyle{definition}
\newtheorem{definition}{Definition}
\theoremstyle{remark}
\newtheorem{remark}{Remark}

\newcommand{\R}{\mathbb{R}}
\newcommand{\Z}{\mathbb{Z}}
\newcommand{\T}{\mathsf{T}}
\newcommand{\code}[1]{\texttt{#1}}
\newcommand{\Lam}{\Lambda}

\title[Two-port substitution and the smallest current in planar networks]{Two-Port Substitution and the Growth\\ of the Smallest Current in Planar Networks}
\author{Filip Majchrzak}
\address{Independent researcher, Poland (SAPIO MORPHEUS project)}
\thanks{ORCID 0009-0007-3475-8900. Version 1.0, 24 November 2026. Archived at Zenodo, \href{https://doi.org/10.5281/zenodo.24497700}{doi:10.5281/zenodo.24497700}. Text licensed CC BY 4.0; code Apache-2.0.}
\date{24 November 2026}

\begin{document}

\begin{abstract}
In a planar network of $q$ unit resistors, feed one unit of current through two terminals on a common face. How
small can a nonzero branch current be? Let $\Lam_q$ be the largest value of $1/|I|$. By the classical
correspondence of Brooks, Smith, Stone and Tutte, $\Lam_q$ is the largest ratio of the longest side to the
smallest square in a rectangle tiled by at most $q$ squares. We show that $c=\lim\Lam_q^{1/q}$ exists and describe
it through two-port networks: every network $X$ with a marked edge has a $2\times2$ matrix $N_X$ of determinant~$1$
such that $c\ge\rho(N_X)^{1/(q-1)}$, where $\rho$ is the spectral radius, and $c$ is the supremum of these numbers.
When $N_X$ is positive, a transposition identity improves this to $c\ge\sigma_{\max}(N_X)^{1/(q-1)}$. With an explicit $23$-edge network we
obtain
\[
 c\;\ge\;1.7194145\ldots,
\]
above the golden ratio $1.618\ldots$ given by the Fibonacci tiling, which is the only lower witness we found in the
literature. The best upper bound known to us is $1.8393$, due to Stoimenow. The substitution lemma, the matrix
form, the iteration and the resulting growth bound are formalised in Lean~4 with Mathlib; the planarity of the
iterated networks, the positivity of the matrix and the hypotheses of the Fekete step are explicit hypotheses
there. The claims about the two explicit networks are checked in exact arithmetic. We do not
claim priority for the question, and we state which parts are classical.
\end{abstract}

\maketitle

\section{Introduction}

Let $H$ be a connected planar multigraph whose edges are unit resistors, and let $s\ne t$ be vertices on a common
face. Feeding one unit of current in at $s$ and out at $t$ produces a current $I_e$ on every edge $e$. Some of
these currents can be exponentially small in the number of edges. The fan (a hub joined to a path) gives the
current $1/\mathrm{Fib}_{2k}$, the reciprocal of a Fibonacci number, with $2k-1$ edges, that is, a decay rate of
the golden ratio per edge.

\begin{definition}
A \emph{gadget} $X=(H,s,t,e)$ consists of a connected planar multigraph $H$ with $q\ge1$ edges (loops and
multiple edges are allowed), distinct vertices
$s,t$ on a common face of some plane embedding, and an edge $e$ with $I_e\ne0$. Let $\Lam_q$ be the largest value
of $1/|I_e|$ over gadgets with at most $q$ edges.
\end{definition}

\subsection*{Squared rectangles}
Brooks, Smith, Stone and Tutte \cite{BSST} showed that a gadget in which every edge carries current is the same
thing as a rectangle $W\times H_0$ tiled by $q$ squares: each edge is a square whose side is its current, the total
current is the width $W$, and the resistance between the terminals is $H_0/W$. Hence $\Lam_q$ is the largest ratio
\[
\frac{\text{longest side of the rectangle}}{\text{side of the smallest square}}
\]
over tilings by at most $q$ squares. (A gadget may have edges without current, which gives fewer squares, and a
tiling in which four squares meet at a point corresponds to several graphs.)

\subsection*{Known bounds}
The literature bounds neighbouring quantities. For a tiling, scale it by the smallest factor that makes all sides
integers. Let $\mathcal P_n$ be the largest value of the longest side of the scaled rectangle over tilings by at
most $n$ squares, let $\mathcal U_n$ be the same over tilings whose scaling contains a unit square, and let $T_n$
be the largest number of spanning trees of a planar multigraph with at most $n$ edges. Then
\[
\mathcal U_n\;\le\;\Lam_n\;\le\;\mathcal P_n\;\le\;T_n .
\]
Conway \cite{Conway} and Kenyon \cite{Kenyon} obtain bounds with base $2$. Fomin \cite{Fomin} gives an elementary
proof of $T_n\le\vartheta^n$ with $\vartheta=1.8637065\ldots$, the largest root of $x^3-x^2-3$, and deduces that
an integer rectangle with a unit square and longest side at least $\vartheta^n$ needs at least $n$ squares.
(See also Bu and Pan \cite{BuPan} for a related determinant bound.) Stoimenow \cite{Stoimenow} proved a bound which, as quoted in \cite{CKPknots}, gives $T_n\le\delta^n$, where
$\delta=1.8393\ldots$ is the real root of $x^3=x^2+x+1$; we rely on that quotation and have not read the original.
Keleti, Lacina, Liu, Liu and Tuir\'an Rangel \cite{Keleti} prove $\mathcal P_n\le4^n$ and note that the Fibonacci
tiling shows that $4$ cannot be replaced by anything below the golden ratio.

For the tiling quantities $\mathcal U_n$ and $\mathcal P_n$ we found in these papers no lower witness above the
golden ratio and no definition of a growth constant. For spanning trees the question is classical: Kenyon
\cite[Problem~2]{Kenyon} asks for the supremum of $\tau(G)^{1/m}$ over planar graphs with $m$ edges and notes that
the square grid gives about $1.79$, so that the supremum lies in $(1.79,2)$.

\subsection*{Results}
\begin{enumerate}
\item The limit $c=\lim_q\Lam_q^{1/q}$ exists and equals $\sup_q\Lam_q^{1/q}$ (Proposition~\ref{prop:limit}).
\item \textbf{Two-port matrix.} Each gadget $X$ has a nonnegative matrix $N_X$ with $\det N_X=1$ such that the
  $k$-fold self-substitution of $X$ has $k(q-1)+1$ edges and its observed current is the reciprocal of the second
  entry of $N_X^k(1,1)^{\T}$ (Lemma~\ref{lem:iter}).
\item \textbf{Characterisation.} $c=\sup_X\rho(N_X)^{1/(q_X-1)}$ (Proposition~\ref{prop:A}). In particular, for
  every rectangle $W\times H_0$ tiled by $q$ squares and every square of side $\ell$ in it,
  $c\ge(\min(W,H_0)/\ell)^{1/(q-1)}$ (Theorem~\ref{thm:B}).
\item \textbf{Transposition.} A role swap followed by planar duality realises $N_X^{\T}$ with the same number of
  edges, so $c\ge\sigma_{\max}(N_X)^{1/(q-1)}$ whenever $N_X$ is positive (Theorem~\ref{thm:C}).
\item \textbf{A bound.} $c\ge1.7194145\ldots$ from an explicit $23$-edge gadget (Section~\ref{sec:numbers}).
\item \textbf{Formalisation.} The substitution lemma, the matrix form, composition, iteration and the growth bound
  are checked in Lean~4 (Section~\ref{sec:lean}).
\end{enumerate}
The constant therefore lies in $[1.7194,\,1.8393]$. Its exact value is open.

\section{The two-port matrix}\label{sec:matrix}

Let $X=(H,s,t,e)$ be a gadget and let $e$ have ends $a,b$. Consider the network $H-e$ with two ports: port~1 is
$(s,t)$ and port~2 is $(a,b)$. Assume first that $e$ is not a bridge, so $H-e$ is connected. Let
\[
S=R_{H-e}(s,t),\qquad P=R_{H-e}(a,b),\qquad Q=\phi(a)-\phi(b),
\]
where $\phi$ is the potential of $H-e$ under the unit $s$--$t$ feed. By reciprocity, $Q$ is also the voltage
between $s$ and $t$ under the unit $a$--$b$ feed.

\begin{lemma}[M\"obius form]\label{lem:mobius}
Give $e$ the resistance $r>0$. Then the resistance between $s$ and $t$ and the current through $e$ under the unit
$s$--$t$ feed are
\[
R_X(r)=\frac{S\,r+(SP-Q^2)}{r+P},\qquad J_X(r)=\frac{Q}{r+P}.
\]
\end{lemma}
\begin{proof}
Let $\phi_1,\phi_2$ be the potentials of $H-e$ under the unit feeds at ports 1 and 2. If $e$ carries current $J$
from $a$ to $b$, the rest of the network sees a unit feed at port~1 and a feed $-J$ at port~2, so its potential is
$\phi_1-J\phi_2$. Ohm's law on $e$ gives $rJ=Q-JP$, hence $J=Q/(r+P)$. The voltage at port~1 is $S-JQ$.
\end{proof}

Let $\gamma=\tau(H-e)$ be the number of spanning trees of $H-e$ and put
\[
M_X=\gamma\begin{pmatrix}S& SP-Q^2\\ 1& P\end{pmatrix}=\begin{pmatrix}\alpha&\beta\\ \gamma&\delta\end{pmatrix},
\qquad \kappa=\gamma\,|Q| ,\qquad N_X=\frac1\kappa M_X .
\]
Then $\det M_X=\kappa^2$ and $\det N_X=1$. By the matrix-tree theorem and Kirchhoff's formulas the entries are
integers with a combinatorial meaning. Let $F(G)$ be the number of spanning forests of $G$ with two components
that separate $s$ from $t$ (so $F(G)=0$ if $s$ and $t$ coincide in $G$). Then
\[
\alpha=F(H-e),\quad \beta=F(H/e),\quad \gamma=\tau(H-e),\quad \delta=\tau(H/e),
\]
and $\kappa=\tau(H)\,|I_e|$ is a positive integer. If $e$ is not an $s$--$t$ edge and $H'=H+f$ for a new edge
$f=st$, then $F(G)=\tau((G+f)/f)$, so $\alpha=\tau(H'/f-e)$, $\beta=\tau(H'/f/e)$, $\gamma=\tau(H'-f-e)$ and
$\delta=\tau(H'-f/e)$. We use this interpretation only in Theorem~\ref{thm:C} and for
exact computation. If $e$ is a bridge, we set $M_X=\bigl(\begin{smallmatrix}\alpha&\beta\\0&\delta\end{smallmatrix}\bigr)$
with the same tree counts; then $\kappa=\alpha=\delta$.

Write $\tau=\tau(H)=\gamma+\delta$, $R=R_H(s,t)$ and $I=I_e$. The row sums of $N_X$ are
\begin{equation}\label{eq:rows}
\frac{\alpha+\beta}{\kappa}=\frac{R}{|I|},\qquad\frac{\gamma+\delta}{\kappa}=\frac1{|I|} .
\end{equation}
For a squared rectangle $W\times H_0$ in which $e$ is a square of side $\ell$, these are $H_0/\ell$ and
$W/\ell$.

\begin{remark}[Classical background]
Up to sign conventions, $N_X$ is the transmission (chain, or $ABCD$) matrix of the reciprocal two-port $H-e$ of
circuit theory: its entries are $S/|Q|$, $(SP-Q^2)/|Q|$, $1/|Q|$ and $P/|Q|$, and $\det N_X=1$ expresses
reciprocity. The cascade rule behind Lemma~\ref{lem:iter} and the matrix identities of Theorem~\ref{thm:C} (port
reversal and duality) are classical in that setting. What we use is that these operations are realised by planar
networks of unit resistors with a controlled number of edges.
\end{remark}

\begin{remark}\label{rem:pos}
$N_X$ is nonnegative. It is positive unless $e$ is a bridge ($\gamma=0$) or $e$ joins $s$ to $t$ ($\beta=0$). In
these two cases $\rho(N_X)=1$, and all bounds below reduce to $c\ge1$.
\end{remark}

\section{Substitution and iteration}

\begin{lemma}[Substitution]\label{lem:subst}
Let $X=(H,s,t,e)$ be a gadget and let $Y$ be a connected network with distinct ports $a',b'$ and resistance $R_Y$
between them. Replace $e$ by $Y$, gluing $a',b'$ to the ends of $e$. Under the unit $s$--$t$ feed, the resistance
between $s$ and $t$ is $R_X(R_Y)$, the current entering $Y$ is $J_X(R_Y)$, and every current inside $Y$ is
$J_X(R_Y)$ times its value under the unit feed of $Y$ at $a',b'$. If $Y$ is planar with $a',b'$ on a common face,
the new network is planar and $s,t$ are still on a common face.
\end{lemma}
\begin{proof}
Scale the unit-feed solution of $Y$ by the current $J$ entering it. Its port voltage is $JR_Y$, so the outer
network cannot distinguish $Y$ from a resistor $R_Y$. For the embedding, draw $Y$ in a thin disk around $e$ with
its ports on the boundary of the disk; faces away from $e$ are unchanged, and a face along $e$ keeps the vertices
it had.
\end{proof}

\begin{lemma}[Iteration]\label{lem:iter}
Let $X$ be a gadget with $q$ edges. Let $G_1=H$, and let $G_{k+1}$ be obtained from $G_k$ by replacing the observed
edge of the innermost copy by a new copy of $(H,s,t)$; observe the edge $e$ of the new copy. Let
$(x_k,y_k)^{\T}=N_X^{\,k}(1,1)^{\T}$. Then $G_k$ has $k(q-1)+1$ edges, its resistance between the outer terminals
is $x_k/y_k$, and its observed current satisfies
\[
\frac1{|I_k|}=y_k .
\]
The same holds for a sequence of different gadgets $X_1,\dots,X_k$ ($X_1$ outermost), with
$(x_k,y_k)^{\T}=N_{X_1}\cdots N_{X_k}(1,1)^{\T}$ and $\sum_i(q_{X_i}-1)+1$ edges.
\end{lemma}
\begin{proof}
A single unit edge has resistance $1=x_0/y_0$. By Lemma~\ref{lem:mobius},
$R_X(x/y)=(\alpha x+\beta y)/(\gamma x+\delta y)$, so the resistance vector is multiplied by $M_X$, and by $N_X$
after scaling. By Lemma~\ref{lem:subst} the observed current is multiplied by
$|J_X(x/y)|=\kappa y/(\gamma x+\delta y)$, which is the ratio of consecutive second coordinates for $N_X$.
\end{proof}

\section{The growth constant}

\begin{proposition}\label{prop:limit}
$\Lam_{p+q}\ge\Lam_p\Lam_q$, and $c=\lim_q\Lam_q^{1/q}=\sup_q\Lam_q^{1/q}$ exists with $1\le c\le2$.
\end{proposition}
\begin{proof}
Take gadgets $X$ and $Y$ attaining $\Lam_p$ and $\Lam_q$. Replace $e$ in $X$ by a path of two unit edges and then
replace the second of them by $Y$. This uses at most $p+q$ edges. The inserted branch has resistance $1+R_Y\ge1$, and
$|J_X(r)|$ is non-increasing in $r$, so by Lemma~\ref{lem:subst} the observed current is at most
$|I_X|\,|I_Y|$ in absolute value and nonzero. The bounds $1\le\Lam_q\le T_q\le2^q$ hold because
$|I|\le1$, $|I|=\kappa/\tau\ge1/\tau$, and a planar multigraph with $q$ edges has at most $2^q$ spanning trees \cite{Kenyon}.
Fekete's lemma \cite{Fekete} gives the limit. (Plain substitution, without the extra series edge, does not
multiply the currents: replacing the observed edge of a triangle, where $1/|I|=3$, by a double edge, where
$1/|I|=2$, gives $1/|I|=5<6$.)
\end{proof}

\begin{theorem}[Free-edge inequality]\label{thm:B}
For every gadget $X$ with $q\ge2$ edges,
\[
c\;\ge\;\rho(N_X)^{1/(q-1)}\;\ge\;\Bigl(\frac{\min(1,R)}{|I|}\Bigr)^{1/(q-1)} .
\]
In particular, for every rectangle $W\times H_0$ tiled by $q\ge2$ squares, each of its graphs and every square of
side $\ell$ in it, $c\ge(\min(W,H_0)/\ell)^{1/(q-1)}$.
\end{theorem}
\begin{proof}
If $N_X$ is not positive, the claim is $c\ge1$ by Remark~\ref{rem:pos}. Otherwise Perron--Frobenius gives $a>0$
with $y_k\ge a\,\rho(N_X)^k$ in Lemma~\ref{lem:iter}. The graph $G_k$ is a gadget with $k(q-1)+1$ edges by
Lemma~\ref{lem:subst}, so $\Lam_{k(q-1)+1}\ge a\rho^k$ and $c\ge\rho^{1/(q-1)}$. The spectral radius of a
nonnegative matrix is at least its least row sum, and the row sums are given by \eqref{eq:rows}.
\end{proof}

Fekete's supremum alone gives $c\ge(\max(W,H_0)/\ell)^{1/q}$. Since $\rho(N_X)\le\max(W,H_0)/\ell$, the gain of
Theorem~\ref{thm:B} is the exponent $1/(q-1)$; the two bounds are incomparable in general, and
Theorem~\ref{thm:B} is stronger for rectangles that are close to squares.

\begin{proposition}[Characterisation]\label{prop:A}
$c=\sup_X\rho(N_X)^{1/(q_X-1)}$, the supremum being over all gadgets.
\end{proposition}
\begin{proof}
``$\ge$'' is Theorem~\ref{thm:B}. For ``$\le$'' we may assume $c>1$. Take a gadget with $q'\le q$ edges and
$1/|I|=\Lam_q$. The conductance between $s$ and $t$ is at most the number of edges, so $R\ge1/q$, and
Theorem~\ref{thm:B} gives $\rho(N_X)\ge\Lam_q/q$, which is at least $1$ for large $q$. Then
$\rho(N_X)^{1/(q'-1)}\ge(\Lam_q/q)^{1/(q-1)}\to c$.
\end{proof}

\begin{theorem}[Transposition]\label{thm:C}
Let $X=(H,s,t,e)$ be a gadget in which $e$ is neither a bridge nor an $s$--$t$ edge, and let
$\Pi=\bigl(\begin{smallmatrix}0&1\\1&0\end{smallmatrix}\bigr)$.
\begin{enumerate}
\item The \emph{role swap} $X^\dagger=(H+f-e;\ \text{terminals } a,b;\ \text{observed edge } f)$, with $f=st$, is a
  gadget with $q$ edges and $N_{X^\dagger}=\Pi N_X^{\T}\Pi$.
\item The \emph{dual gadget} $X^\ast$, obtained from the plane dual of $H+f$ by deleting $f^\ast$, with the ends of
  $f^\ast$ as terminals and $e^\ast$ observed, is a gadget with $q$ edges and $N_{X^\ast}=\Pi N_X\Pi$.
\item Hence $(X^\dagger)^\ast$ realises $N_X^{\T}$, and
\[
c\;\ge\;\rho\bigl(N_XN_X^{\T}\bigr)^{1/(2(q-1))}=\sigma_{\max}(N_X)^{1/(q-1)}\;\ge\;\rho(N_X)^{1/(q-1)} .
\]
\end{enumerate}
\end{theorem}
\begin{proof}
(1) The graph $H+f$ is planar because $s$ and $t$ share a face. The graph $H+f-e$ is connected, and $a,b$ lie on
the face obtained by merging the two faces along $e$. Its
two-port network after removing the observed edge is again $H-e$, with the two ports exchanged. So $S$ and $P$
are exchanged and $Q$ is unchanged, which turns $\bigl(\begin{smallmatrix}S&D\\1&P\end{smallmatrix}\bigr)$,
$D=SP-Q^2$, into $\bigl(\begin{smallmatrix}P&D\\1&S\end{smallmatrix}\bigr)=\Pi
\bigl(\begin{smallmatrix}S&D\\1&P\end{smallmatrix}\bigr)^{\!\T}\Pi$; the factor $\gamma$ and $\kappa$ are the
same.

(2) Planar duality exchanges deletion and contraction and preserves the number of spanning trees. Applied to the
four tree counts of $H'$ in Section~\ref{sec:matrix} it gives
$(\alpha^\ast,\beta^\ast,\gamma^\ast,\delta^\ast)=(\delta,\gamma,\beta,\alpha)$, that is,
$M_{X^\ast}=\Pi M_X\Pi$, and $\kappa^\ast=\kappa$ because $\det$ is unchanged. Here $e^\ast$ is not a bridge of the
dual gadget, because $\gamma^\ast=\beta>0$, and it does not join its terminals, because $\beta^\ast=\gamma>0$. The
terminals of $X^\ast$ are joined by $f^\ast$ in a planar graph, so they are on a common face.

(3) Alternate $X$ and $(X^\dagger)^\ast$ in Lemma~\ref{lem:iter}: one period has $2(q-1)$ edges and the matrix
$N_XN_X^{\T}$, which is positive, and the argument of Theorem~\ref{thm:B} applies. The last inequality is
$\sigma_{\max}\ge\rho$.
\end{proof}

Note that $\operatorname{tr}(N_XN_X^{\T})$ is the sum of the squares of the entries of $N_X$.

\section{Explicit gadgets and the bound}\label{sec:numbers}

The numbers for $H_7$ and $X_{23}$ below were computed in exact rational arithmetic, by two separately written
programs. The catalogue table is a floating-point computation.

\subsection*{An $11$-edge example}
Let $H_7$ have vertices $0,\dots,6$ and edges
$04,05,06,13,15,16,23,24,26,36,45$, with $s=3$, $t=4$ and $e=05$. Then $\tau=209$, $I_e=-1/209$ and
\[
M=\begin{pmatrix}83&94\\98&111\end{pmatrix},\qquad\kappa=1,\qquad \operatorname{tr}M=194 .
\]
The iterated graphs have $10k+1$ edges and currents $1/209,\ 1/40545,\ 1/7865521,\dots$, with
$y_{k+2}=194y_{k+1}-y_k$, and $\rho(M)=97+56\sqrt3=(2+\sqrt3)^4$. This gives
$c\ge(2+\sqrt3)^{2/5}=1.6934758\ldots$ The corresponding tiling, $209\times177$, is in the catalogue of simple
perfect squared rectangles of order~$11$ \cite{Anderson}.

\subsection*{Catalogue scan}
We evaluated $\rho(N_X)^{1/(q-1)}$ for every simple perfect squared rectangle of orders $11$ to $21$ in the
catalogue \cite{Anderson} (about $6.9$ million tilings), with every square as the observed edge; by
Theorem~\ref{thm:C}(2) the two orientations of a tiling give the same value. The ranking used floating-point
arithmetic; the best gadgets of orders $11$, $19$, $20$ and $21$ were then
verified exactly.
\begin{center}
\small
\begin{tabular}{lccccccccccc}
\toprule
order&11&12&13&14&15&16&17&18&19&20&21\\
\midrule
best rate&1.6935&1.6885&1.6898&1.6907&1.7022&1.7055&1.7062&1.7096&1.7141&1.7144&1.7150\\
\bottomrule
\end{tabular}
\end{center}

\subsection*{A $23$-edge gadget}
A bounded computer search beyond the catalogue, seeded with the best order-$20$ gadgets, produced the following
gadget $X_{23}$ on $13$ vertices with $23$ edges:
\begin{gather*}
02,\ 01,\ 12,\ 16,\ 24,\ 45,\ 47,\ 59,\ 58,\ 67,\ 78,\ 89,\ 17,\ (10)8,\ (10)6,\ (10)0,\ (11)9,\ (10)(11),\\
30,\ 32,\ (12)4,\ (12)(11),\ 3(12);\qquad s=12,\quad t=5,\quad e=(10)6 .
\end{gather*}
It is simple and planar, $s,t$ are on a common face, and
\[
\tau=160192,\qquad |I_e|=\frac1{160192},\qquad M=\begin{pmatrix}63678&75857\\73105&87087\end{pmatrix},
\qquad\kappa=1,\qquad\det M=1 .
\]
It corresponds to a simple perfect squared rectangle $160192\times139535$ of order $23$ with a unit square. The
search is heuristic; at order $21$ it returned the optimum of the complete catalogue.

By Theorem~\ref{thm:B}, $c\ge\rho(M)^{1/22}=1.7194015\ldots$ By Theorem~\ref{thm:C}, with
$\operatorname{tr}(MM^{\T})=22\,737\,658\,727$,
\[
c\;\ge\;\Bigl(\tfrac12\bigl(22737658727+\sqrt{22737658727^2-4}\bigr)\Bigr)^{1/44}=1.7194145008\ldots
\]
We checked the construction directly. The gadget $(X_{23}^\dagger)^\ast$ is simple and planar, has $23$ edges,
terminals on a common face, and matrix $M^{\T}$. The networks $X_{23}\circ(X_{23}^\dagger)^\ast$ ($45$ edges) and
its square ($89$ edges) are simple and planar with terminals on a common face, and their observed currents are
exactly $1/24\,189\,825\,343$ and $1/550\,019\,993\,314\,869\,718\,360$, as predicted by the matrix products;
both denominators equal the numbers of spanning trees. The bound holds a fortiori for $\mathcal P_n$ and $T_n$.
Since these gadgets have $\kappa=1$, it also holds for $\mathcal U_n$.

\section{Lean formalisation}\label{sec:lean}

The general method is formalised in Lean~4 (\code{leanprover/lean4:v4.37.0-rc1}) with Mathlib (commit
\code{53dad9ed}), in a package \code{TwoPort} of about $1350$ lines. It imports the Fekete module of the
development \cite{M1} and is distributed in the same repository. Networks are described by potentials and
Kirchhoff's law with real resistances, positivity being assumed where it is needed; chains and iteration are
stated for unit resistances. No Laplacian determinants and no matrix-tree theorem are used.
\begin{center}
\small
\begin{tabularx}{\linewidth}{>{\ttfamily}l X}
\toprule
\textrm{statement}&content\\
\midrule
Net.subst\_isSol& Lemma~\ref{lem:subst}: outer currents unchanged, inner currents multiplied by the branch current\\
Port.resp\_of\_data, det\_mobM& Lemma~\ref{lem:mobius} and $\det=Q^2$; also $S>0$, $P\ge0$, $Q^2\le SP$\\
Port.hasMob\_comp, chain\_current& composition is the matrix product; chains of arbitrary lists of gadgets\\
Gadget.ofData& the normalised matrix $N_X$, with $\det N_X=1$, for every unit-resistance two-port given with its
open-circuit potentials and $Q\ne0$\\
iterate\_current& Lemma~\ref{lem:iter}: $(q-1)k+1$ edges and $|I_k|=1/y_k$\\
min\_rowSum\_le\_lam& $\lambda\ge$ least row sum for a positive matrix of determinant $1$, where
$\lambda=\tfrac12(\operatorname{tr}+\sqrt{\operatorname{tr}^2-4})$ is the larger root of its characteristic
polynomial\\
growth\_of\_gadget, growth\_fekete& $\lambda(N_X)^{1/(q-1)}\le c$, under the hypotheses below\\
H7.gadget\_M, H7.q\_rec, H7.lam\_eq\_pow& the example $H_7$: matrix, recurrence $194$, $\lambda=(2+\sqrt3)^4$\\
\bottomrule
\end{tabularx}
\end{center}
\textbf{Hypotheses that are not formalised.} (i) That the iterated networks are planar with terminals on a common
face: in Lean this is an explicit hypothesis saying that $\Lam$ dominates the inverse observed current of each
iterated network. (ii) The hypotheses of the Fekete step (supermultiplicativity, $\Lam\ge1$, an exponential upper
bound), as in \cite{M1}. (iii) The interpretation of the matrix entries as spanning-tree counts. (iv) The
positivity of the entries of $N_X$, which is a hypothesis of the growth theorems in Lean and is checked for $H_7$;
Remark~\ref{rem:pos} is not formalised. (v) The existence of the open-circuit potentials of $H-e$, which are part
of the input data of a two-port in Lean.
\textbf{Not formalised at all:} Theorem~\ref{thm:C}; the direction ``$\le$'' of Proposition~\ref{prop:A}; the
identification of the row sums of $N_X$ with $R/|I|$ and $1/|I|$, and hence the statement of
Theorem~\ref{thm:B} about squared rectangles; the identification of the graphs $G_k$ of Lemma~\ref{lem:iter} with
the chains used in Lean; that $\lambda$ is the spectral radius; and the numerical results of
Section~\ref{sec:numbers} other than $H_7$.

The package builds without \code{sorry}. The $21$ theorems listed in \code{scripts/PrintAxioms.lean}, and the
statements of the table that are not in that list (checked separately), depend only on the axioms
\code{propext}, \code{Classical.choice} and \code{Quot.sound}. Every module is replayed by \code{leanchecker}.
The file \code{growth-constant/VERIFICATION\_CLEANROOM.txt} of the archive records the outputs of a build from an
empty directory with a freshly downloaded Mathlib.

\section{Status of the claims and open questions}

\begin{center}
\small
\begin{tabularx}{\linewidth}{l X}
\toprule
classical& Kirchhoff's laws; the matrix-tree theorem; transmission matrices of two-ports, their cascade rule,
port reversal and duality; planar duality; the correspondence with squared rectangles \cite{BSST}; Fekete's
lemma; the Fibonacci tiling\\
prior work& upper bounds $2$ \cite{Conway,Kenyon}, $1.8638$ \cite{Fomin} and $1.8393$ \cite{Stoimenow};
$\mathcal P_n\le4^n$ \cite{Keleti}; Kenyon's question on the spanning-tree constant; the catalogue of squared
rectangles \cite{Anderson}, which contains the tiling of $H_7$\\
possibly new& the growth constant $c$ of the ratio (longest side)/(smallest square) and its existence; the use of
two-port matrices with the exponent $1/(q-1)$ (Theorem~\ref{thm:B}, Proposition~\ref{prop:A}); the realisation of
the transposed matrix by a planar gadget with the same number of edges (Theorem~\ref{thm:C}); a lower bound above
the golden ratio; the Lean development\\
\bottomrule
\end{tabularx}
\end{center}
Our literature search was limited (web searches and the full texts of \cite{Kenyon,Fomin,Keleti}). The method is
elementary, and we would be grateful for references to earlier occurrences. We do not claim priority.

\textbf{Open questions.} (1) The value of $c$; we have $1.7194\le c\le1.8393$. (2) Let
$c_\tau=\lim T_n^{1/n}$. Kenyon asked for this constant \cite[Problem~2]{Kenyon}; the conjecture
$c_\tau=e^{2C/\pi}=1.7916\ldots$, with $C$ Catalan's constant, is attributed to him in \cite{CKPknots} and is
also stated by Fomin \cite{Fomin}. Is $c=c_\tau$? Heuristically, networks with nearly the maximal number of
spanning trees resemble lattices, in which currents decay only polynomially, so we expect $c<c_\tau$; we have no
proof. (3) For terminals that are not on a common face the
correspondence with tilings is lost. Is the growth constant the same?

\section*{Author contributions and use of AI}
This work was carried out with substantial assistance from AI systems working under the author's direction. The
SAPIO MORPHEUS project uses models from OpenAI and Anthropic.
\begin{itemize}
\item The example $H_7$ and its recursion were first found by an AI agent of the project based on an OpenAI model.
\item The verification, the two-port method, the theorems and their proofs, the catalogue analysis, the search for
  $X_{23}$, the Lean formalisation and this manuscript were produced by AI agents based on Anthropic's Claude.
\item An independent AI referee checked the proofs and recomputed the numerical claims with separate code; it
  found Theorem~\ref{thm:C}.
\item The author selected the problem, directed the work and made the publication decisions, and takes
  responsibility for the content.
\end{itemize}
In line with common publication policies, AI systems are not listed as authors.

\begin{thebibliography}{99}
\bibitem{Anderson} S.~Anderson, \emph{Catalogues of simple perfect squared rectangles}, \url{www.squaring.net}
  (accessed 24 November 2026).
\bibitem{BSST} R.~L.~Brooks, C.~A.~B.~Smith, A.~H.~Stone, W.~T.~Tutte, \emph{The dissection of rectangles into
  squares}, Duke Math.\ J.\ \textbf{7} (1940), 312--340.
\bibitem{BuPan} A.~Bu, Y.~Pan, \emph{Planarity via spanning tree number: a linear-algebraic criterion}, SIAM J.\
  Discrete Math.\ \textbf{39} (2025), 728--751.
\bibitem{CKPknots} A.~Champanerkar, I.~Kofman, J.~S.~Purcell, \emph{Geometrically and diagrammatically maximal
  knots}, arXiv:1411.7915v4 (2018).
\bibitem{Conway} J.~H.~Conway, \emph{Mrs Perkins's quilt}, Proc.\ Cambridge Philos.\ Soc.\ \textbf{60} (1964),
  363--368.
\bibitem{Fekete} M.~Fekete, \emph{\"Uber die Verteilung der Wurzeln bei gewissen algebraischen Gleichungen mit
  ganzzahligen Koeffizienten}, Math.\ Z.\ \textbf{17} (1923), 228--249.
\bibitem{Fomin} D.~Fomin, \emph{Elementary proof for the bounds of the complexity of a planar multigraph and the
  size of a prime rectangular squaring}, arXiv:2103.10523 (2021).
\bibitem{Keleti} T.~Keleti, S.~Lacina, C.~Liu, M.~Liu, J.~R.~Tuir\'an Rangel, \emph{Tiling of rectangles with
  squares and related problems via Diophantine approximation}, Discrete Math.\ \textbf{346} (2023), 113442.
\bibitem{Kenyon} R.~Kenyon, \emph{Tiling a rectangle with the fewest squares}, J.\ Combin.\ Theory Ser.\ A
  \textbf{76} (1996), 272--291.
\bibitem{M1} F.~Majchrzak, \emph{Machine-checked flow-certificate thresholds for planar graphs using
  rotation-system certificates}, Zenodo (2026), doi:10.5281/zenodo.24496271.
\bibitem{Stoimenow} A.~Stoimenow, \emph{Maximal determinant knots}, Tokyo J.\ Math.\ \textbf{30} (2007),
  73--97.
\end{thebibliography}

\end{document}
